% Transcription — fonds Alexandre Grothendieck, shelfmark 156-5, batch 3 (pages 41-48).
% Established from the facsimile archives/batches/156-5/batch-03.pdf (a mirror of
% https://grothendieck.umontpellier.fr/156-5.pdf; no cover sheet, PDF page k =
% folder page 40+k), per the skill transcribe-grothendieck. The folder is 48
% pages; this is its last batch, eight pages.
%
% All eight leaves are mathematics and all are transcribed; there is no cover
% and no unrelated verso. His own page numbers 41-48 stand at the head of each
% leaf and coincide with the archivists'. Page 41 opens in the middle of an
% enumeration (item c) of a proposition whose statement lies in batch 2; the
% notation is taken as it stands here.
%
% Content. The setting: a set S with a topology (or ultratopology, i.e. a
% preorder) tau, a closed part T with a topology tau'_0 on T, U = S \ T, and
% the topology tau' on S obtained by gluing: its closed sets are the pairs
% (B, A_U), B closed in T for tau'_0, A_U closed in U for tau|U, with
% cl_tau(A_U) ∩ T ⊂ B (p. 41, c). d) describes the preorder of tau' in the
% ultratopological case, case by case. e) passes to the associated ordered
% sets phi : I' -> I, J = image of T, J' = phi^{-1}(J), and characterises the
% topology on I' by four conditions (i)-(iv) (p. 42); a slanted margin note on
% p. 42 already says « C'est faux » and points to p. 48. Pages 42-44: proof of
% c) and d) (tau' is the coarsest topology of the class, the closure formula
% cl_tau'(X) = cl_tau(X ∩ U) ∪ cl_tau'_0(X ∩ T), written in the margin of
% p. 44). Page 45: a corollary (tau is final for S -> S, T -> S), with a long
% middle passage cancelled by two diagonal strokes. Page 46: ultratopological
% corollaries on phi, phi^* : surjectivity, « homotopiquement propre » with
% « faiblement contractiles » fibres (weak homotopy equivalence, « kif-kif »
% for phi^*), an argument by greatest element closed with « OK ! ». Page 47:
% he notices e) is not proved (« Mais je n'ai pas démontré … »; margin
% « c'est faux !! »), sets out the conditions (i)-(iv) for a tau'' and
% reduces to cl_tau''(A) = cl_tau(A) for A ⊂ U. Page 48: « ?? », then « C'est
% ultra faux », a counterexample on a three-element I' over a two-element I
% (drawn as two ovals with points and arrows; described, not redrawn), and a
% ruled-off conclusion: the condition to impose is that closures of parts of
% U agree for tau and tau'.
%
% Reading hazards fixed once here. tau with an overbar (\overline{\tau}) is the
% topology on the ordered set I; with a tilde (\widetilde{\tau}) the inverse
% image topology on I'; tau'_0 is written τ with a prime and a small 0 below.
% \underline{\mathcal{C}} is his underlined curly class of topologies (the
% letter is looped like a τ and could be read \underline{\tau}). His
% underlining in prose is set in \emph. X_U abbreviates X ∩ U, X_B the closed
% part B of X. Overlined sets with a superscript (τ) are closures for that
% topology, set as \overline{X}^{(\tau)}.
%
% Readings flagged and not settled: most of the connective prose of the last
% paragraph of p. 41; the slanted notes of p. 42 (top-left and bottom-left) and
% the circled insertion and slanted block of p. 42's middle; the corollary
% statement of p. 44; the opening of p. 45 and the letter T/S in its
% corollary; the first line of p. 47. Variances on the page and not repaired:
% p. 43 writes τ'|U ≠ τ|U where the argument needs =; p. 48 writes ∩ in the
% closure formula that p. 44's margin writes with ∪; p. 42 (iii) writes
% U = S \ T where U = I \ J would be expected.
%
% Pass: Opus 5.5 (claude-opus-5-5), 2026-09-24 — first pass, unchecked against the pages by a human.
\documentclass[11pt,a4paper]{article}
% Le préambule commun à toutes les transcriptions du fonds Grothendieck.
%
% Il tient en une idée : **l'incertitude se compose, elle ne se résout pas.**
% Une transcription qui lisse un mot douteux a détruit la seule chose pour
% laquelle on la faisait — un lecteur doit pouvoir savoir, à chaque endroit,
% ce qui était sur le feuillet et ce qui a été ajouté. Les six macros qui
% suivent sont donc l'appareil critique, et non de la décoration.
%
% Les mêmes commandes sont reconnues par `scripts/render.mjs`, qui produit la
% vue de lecture du site. Une macro ajoutée ici doit l'être aussi là-bas,
% faute de quoi le rendu échouera bruyamment — ce qui est voulu : un rendu
% silencieusement faux est pire qu'un rendu absent.

\ProvidesPackage{grothendieck}[2026/08/08 Transcriptions du fonds Grothendieck]

\usepackage{amsmath,amssymb,amsthm}
\usepackage{mathrsfs}
\usepackage{stmaryrd}
% Les diagrammes commutatifs sont partout dans ce fonds : c'est la langue
% même de la géométrie algébrique de Grothendieck, et une transcription qui
% les rendrait en prose ne transcrirait rien.
\usepackage{tikz-cd}
% Français par défaut — c'est la langue des feuillets — mais l'anglais est
% chargé aussi : la traduction et le résumé sont des documents anglais, et la
% typographie française (espaces avant les deux-points) y serait une faute.
% Ces documents basculent par \selectlanguage{english} en tête de corps.
\usepackage[english,french]{babel}
\usepackage{fontspec}
\usepackage{geometry}
\geometry{a4paper,margin=25mm,marginparwidth=28mm}
\usepackage{marginnote}
\usepackage{xcolor}
% ulem, not soul. `soul`'s \st and \ul are built on a character-by-character
% scan that XeLaTeX's font machinery breaks: the document fails with an
% undefined control sequence rather than degrading. ulem does the same two jobs
% and is Unicode-engine safe. `normalem` stops it hijacking \emph, which the
% transcriptions use for Grothendieck's own emphasis.
\usepackage[normalem]{ulem}
\usepackage{hyperref}
\hypersetup{colorlinks=true,linkcolor=black,urlcolor=black,pdfborder={0 0 0}}

% --- Métadonnées du lot -----------------------------------------------------
% Reprises telles quelles par le script de rendu, qui les affiche en tête de
% la vue de lecture. Elles fixent ce que le fichier prétend couvrir : sans
% elles, une transcription flotte sans point d'attache dans un fonds de seize
% mille feuillets.

\newcommand{\folder}[1]{\gdef\grothFolder{#1}}
\newcommand{\batch}[1]{\gdef\grothBatch{#1}}
\newcommand{\leaves}[2]{\gdef\grothFirst{#1}\gdef\grothLast{#2}}
\newcommand{\foldertitle}[1]{\gdef\grothFoldertitle{#1}}
\gdef\grothFolder{?}\gdef\grothBatch{?}\gdef\grothFirst{?}\gdef\grothLast{?}\gdef\grothFoldertitle{}

% --- Le résumé --------------------------------------------------------------
%
% Il ouvre la lecture modernisée, sous le titre « Résumé », et il est composé
% en petit. Ce n'est pas un ornement : c'est le seul passage du document qui
% ne soit pas des mathématiques, et un lecteur doit pouvoir le reconnaître
% avant de l'avoir lu — puis le sauter s'il connaît déjà le sujet, ou s'y
% arrêter s'il ne le connaît pas. Le filet qui le suit marque la reprise du
% corps.

\newenvironment{resume}
  {\small\setlength{\parskip}{0.45em}}
  {\par\medskip\hrule\medskip}

% --- Mots-clés ---------------------------------------------------------------
%
% En anglais, en clôture du résumé de la lecture modernisée : le vocabulaire
% moderne sous lequel chercher ce que les feuillets construisent. Le manifeste
% du site (`npm run manifest`) les extrait pour étiqueter la cote entière —
% cette ligne est la source unique des tags, il n'y a pas de fichier à côté.
\newcommand{\keywords}[1]{\par\smallskip\noindent
  {\footnotesize\textsc{Keywords} --- \textit{#1}}\par}

% --- L'appareil critique ----------------------------------------------------

% Le numéro de feuillet, en marge. C'est l'ancre qui permet de régler un
% désaccord sur une formule en regardant *un* feuillet plutôt qu'une chemise
% de six cents. Le site s'en sert aussi pour tourner les pages du fac-similé
% quand on fait défiler la transcription.
\newcommand{\leaf}[1]{%
  \marginnote{\footnotesize\color{gray!70}\textsf{#1}}%
  \label{leaf:#1}\ignorespaces}

% Illisible : on ne devine pas. Un mot inventé qui se lit comme les autres est
% le pire résultat possible.
\newcommand{\ill}{\textcolor{red!60!black}{[\,\ldots\,]}}

% Lecture douteuse : le mot est donné, et signalé comme incertain.
\newcommand{\uncertain}[1]{\uline{#1}}

% Ajout de l'éditeur — une lettre manquante, un mot sous-entendu.
\newcommand{\add}[1]{\textcolor{blue!50!black}{[#1]}}

% Biffé par Grothendieck lui-même. Ce qu'il a rayé fait partie du document :
% une rature dit souvent où la pensée a changé de direction.
\newcommand{\struck}[1]{\sout{#1}}

% Note du transcripteur, sur l'état matériel du feuillet.
\newcommand{\note}[1]{\par\smallskip{\small\color{gray!60!black}\textit{[#1]}}\par\smallskip}

% Note marginale de Grothendieck, distincte du corps du texte.
\newcommand{\marginal}[1]{\par\smallskip\noindent\hspace{1em}%
  {\small\color{gray!60!black}#1}\par\smallskip}

% --- Titre ------------------------------------------------------------------

\AtBeginDocument{%
  \begin{center}
    {\large\bfseries Fonds Alexandre Grothendieck --- cote n\textsuperscript{o}~\grothFolder}\\[2pt]
    {\itshape\grothFoldertitle}\\[4pt]
    {\small feuillets \grothFirst--\grothLast\ (lot \grothBatch)}\\[2pt]
    {\footnotesize\color{gray!60!black}Transcription établie d'après le fac-similé numérisé
     par l'Université de Montpellier.\\
     Les lectures incertaines sont soulignées, les illisibles notées [\,\ldots\,],
     les ajouts entre crochets.}
  \end{center}
  \vspace{1em}}

\endinput


\folder{156-5}
\batch{3}
\pages{41}{48}
\dating{1986}
\watermark{Édition de démonstration}
\foldertitle{[Chapitre] V. Algèbre des figures : notes manuscrites (14/06/1986)}

\begin{document}
\page{41}
\note{la page s'ouvre sur l'item c) d'un énoncé dont le début est à la page précédente}

c) L'ens. des fermés pour $\tau'$ est en corr. biunivoque avec l'ens. des couples
\[
  (B, A_{\mathcal{U}}) \qquad
  \begin{array}{l}
    B \subset T \text{ fermé pour } \tau'_0 \\
    A_{\mathcal{U}} \text{ fermé de } \mathcal{U} = S \smallsetminus T \text{ pour } \tau|(S \smallsetminus T)
  \end{array}
\]
tels que
\[
  \overline{A_{\mathcal{U}}} \cap T \subset B \qquad (\text{adh. prise pour } \tau)
\]
\note{sous « (adh. prise pour $\tau$) », un symbole biffé, peut-être $\tau$}

d) Si $\tau'_0$ et $\tau$ ultratopologiques, $\tau'$ aussi. La relation de préordre associée \struck{se définit} se décrit ainsi. Pour $x, y \in S$, on a $x \leq_{\tau'} y$ ssi
\[
  \begin{cases}
    \text{ou bien } x, y \in T, & x \leq_{\uncertain{\tau'_0}} y \\
    \text{ou bien } x, y \in \mathcal{U} = S \smallsetminus T, & x \leq_{\tau} y \\
    \text{ou bien } x \in T,\ y \in \mathcal{U}, & x \leq_{\tau} y
  \end{cases}
\]
\note{l'indice du premier cas est écrit $\tau'$ suivi d'un petit signe, lu $\tau'_0$}
\marginal{\uncertain{En} d'autres termes, si $y \in \mathcal{U}$, on a $x \leq_{\tau'} y$ ssi $x \leq_{\tau} y$.}

\uncertain{Sous} les conditions de d)\note{ces mots sont écrits au-dessus de la ligne et rattachés à « e) » par un trait}

e) Considérons les ens. ordonnés $I$, $I'$ associés à $\overline{\tau}$, $\overline{\tau}'$\note{les deux lettres sont barrées d'un trait horizontal, lues $\tau$ surligné}, et la partie \add{fermée}\note{« fermée » au-dessus de la ligne} $J \subset I$ correspondant à la partie fermée $T$ de $(S, \tau)$, enfin l'application $\varphi : I' \to I$ canonique.

\begin{tikzcd}
\varphi^{-1}(J) = J' \arrow[r, hook] \arrow[d] & I' \arrow[d, "\varphi"] & \overline{\tau}' \arrow[d] \\
J \arrow[r, hook] & I & \overline{\tau}
\end{tikzcd}

\note{la flèche $\overline{\tau}' \to \overline{\tau}$ est surchargée d'un gribouillis}

\emph{Mais}, \struck{\ill{}} soit $\overline{\tau}'$ \uncertain{est} la top. de $\tau'$ déduite de $\tau'$ par passage au quo\ill{} (qui \uncertain{lient} les \ill{}, \uncertain{associée} \ill{} relation \emph{d'ordre} sur $I'$ \uncertain{associée} au \uncertain{préordre}), $\widetilde{\tau}$ \ill{} $I'$ image inverse de $\overline{\tau}$. Alors \uncertain{est plus fine que} $\widetilde{\tau}$\note{ces mots sont écrits au-dessus de la ligne, après « Alors »} \ill{} la topologie sur $I'$ qui \ill{} \uncertain{induit} \ill{} \ill{} des top. sur $I'$ qui \ill{} topologie $\overline{\tau}|J'$. Elle est

\page{42}
\marginal{C'est faux, \uncertain{sauf} (iii), (iv) \ill{} \ill{} \ill{} \ill{} \ill{} \ill{} \ill{} ! \ill{} — i.e. $\overline{\tau}' \in \mathcal{C}$}
\note{note oblique dans l'angle supérieur gauche, séparée du texte par un trait ; au-dessus, un renvoi \uncertain{« v. 48 »}}

caractérisée par les conditions suivantes :
\begin{itemize}
  \item[(i)] $\overline{\tau}'$ est plus fine que $\widetilde{\tau}$, et induit sur $J'$ $\overline{\tau}|J'$
  \item[(ii)] $\overline{\tau}'|\mathcal{U}' = I' \smallsetminus J'$ coïncide avec $\widetilde{\tau}|\mathcal{U}'$
  \item[(iii)] $\varphi|\mathcal{U}'$ induit une \struck{bijection} \uncertain{iso. ord.} $\mathcal{U}' \to \mathcal{U} = S \smallsetminus T$
  \item[(iv)] La topologie de $I$ est quotient \uncertain{fermé} pour les applications $I' \to I$, $J \to I$ ($I'$ muni \struck{\ill{}} de $\overline{\tau}'$, $J$ muni de $\overline{\tau}|J$)
\end{itemize}
\marginal{(Ce qui \uncertain{inclut} \ill{} \ill{} trivial pour \ill{})}
\marginal{\struck{Au\ill{}}}

(i), (ii), \uncertain{obtenu} en prenant la topologie \uncertain{somme} sur \struck{\ill{}} $I' = J' \cup \mathcal{U}'$ (disj.), \ill{} $I$ lui-même est une somme \ill{} ($J \sqcup \mathcal{U}$, $\mathcal{U} = I \smallsetminus J$).

\marginal{C'est une condition que j'avais d'abord \uncertain{oubliée} \ill{} à la page 31 !}

\emph{Démonstration}

\struck{a) et b).} (a b c) Il est clair que les \ill{} de c) sont \uncertain{les fermés} \ill{} \ill{} b), et que \ill{} est \uncertain{unique}.
\note{dans la marge droite, un paragraphe oblique entouré d'une boucle : « en fait, cela \uncertain{implique} \ill{} \ill{} \ill{} \ill{} \uncertain{ultratop.} » ; puis un bloc oblique, dont le premier mot est biffé : « \ill{} dire que $\overline{\tau}'$, \uncertain{sur} $S$ est \ill{} \uncertain{par} (i) (ii) (iii) (iv) \ill{} revient à dire que \ill{} top. $\tau'$ \ill{} $S$ \ill{} $\tau$ sur $S$ et $\tau'$ sur $T$ ($\uncertain{\mathcal{U}} \supset S$) \ill{} » ; il n'est lu qu'en partie}

\ill{} \ill{} qui satisfait 1°), 2°), 3°) ; que les \uncertain{conséquences} \ill{}. Il reste à voir, pour \uncertain{un} (a, b, c), que les \uncertain{ensembles} \uncertain{ainsi} formés \uncertain{sont} une topologie, et qu'elle \struck{les propr.} \ill{} dans $\underline{\mathcal{C}}$, il est \uncertain{la} moins fine des top. $\in \underline{\mathcal{C}}$.

1) C'est \uncertain{l'ens. des} fermés d'une topologie. Stable par intersections quelconques : clair ; stable par réunions finies : \uncertain{id.} C'est \uncertain{=} stable par réunions quelconques s'il en est ainsi pour $(T, \tau'_0)$, $(\struck{\ill{}}\,\uncertain{\mathcal{U}}, \tau)$, et c'est le \uncertain{préordre} associé de d).

\page{43}
Cette topologie est plus fine que $\uncertain{\tau}$, i.e. si $X$ est fermé dans $\tau$, il satisfait aux conditions : clair. Il induit $\tau'_0$, i.e. si $X \subset T$, il est fermé \uncertain{ssi} il est fermé pour $\tau'_0$ : clair. \emph{Donc} $\tau' \in \underline{\mathcal{C}}$. Prouvons que c'est un plus grand \uncertain{élément} (\uncertain{top. la moins} fine), i.e. que pour tout $\tau'' \in \underline{\mathcal{C}}$, les fermés pour $\tau'$ sont des fermés pour $\tau''$.
\[
  X = B \cup \overline{X \cap \mathcal{U}}^{(\tau)}
\]
En effet, \ill{} \ill{}.

$B$ étant fermé dans $T$ (pour $\tau'_0$), \struck{\ill{}} \uncertain{fermé} pour $\tau''|T$\note{« pour $\tau''|T$ » est écrit au-dessus de la ligne, en remplacement d'un mot biffé}, \emph{fermé} dans $S$ pour $\tau''$\note{au-dessus, « \uncertain{effectivement} »}, donc pour $\tau''$ ;
\marginal{(car \ill{} \ill{} \ill{} \ill{} $T$ fermé dans $S$)}
$\overline{X \cap \mathcal{U}}^{(\tau)}$ est fermé dans $S$ pour $\tau$, \uncertain{l'est} pour $\tau''$, donc
\[
  X = B \cup \overline{X \cap \mathcal{U}}^{(\tau)}
\]
\uncertain{l'est aussi}. \uncertain{Cqfd}.

Prouvons a) b) c). (\uncertain{C'est déjà} \uncertain{fait} \ill{})

Prouvons d), en \uncertain{notant} \ill{} que
\[
  x \leq_{\tau'} y \iff \overline{x}^{(\tau')} \subset \overline{y}^{(\tau')} \iff x \in \overline{y}^{(\tau')} .
\]
Comme $\tau'|\mathcal{U} \neq \tau|\mathcal{U}$, et $\tau'|T = \tau'_0$,
\note{le signe est lu $\neq$ ; l'argument demande $=$, cf. (ii) p. 42 et (iii) p. 47}
les cas $x, y \in T$ et $x, y \in \mathcal{U}$ sont déjà clairs. Reste le cas où \ill{} $x \in \mathcal{U}$, $y \in T$, \ill{} \uncertain{à} \uncertain{prouver}. Le premier cas ne peut \uncertain{se produire}, car si $y \in T$, $\overline{y}^{(\tau')} \subset T$ ($T$ fermé dans $S$ pour $\tau'$) donc $x \notin \overline{y}^{(\tau')}$ \struck{puisque}
si $x \in \mathcal{U}$. \uncertain{Donc} reste le cas [$x \in T$,] $y \in \mathcal{U}$, prouvons $x \leq_{\tau'} y \iff x \leq_{\tau} y$,
\marginal{\textbf{NB} l'autre cas \ill{} \ill{} \ill{} $x \in T$}

\page{44}
donc il faut voir que l'adhérence de $y$ pour $\tau'$, et pour $\tau$, est pareille (a priori \uncertain{elle} serait plus petite). \struck{Soit}

\struck{En effet \ill{} $X$ \ill{} \ill{} \ill{} de $X$, tel que}

\marginal{Cela donne l'adhérence d'une partie \uncertain{générale} \struck{$X$} de $S$, pour $\tau'$, en termes de $\tau$ et $\tau'_0$ : $\overline{X}^{(\tau')} = \left(\overline{X \cap \mathcal{U}}^{(\tau)}\right) \cup \overline{(X \cap T)}^{(\tau'_0)}$}
\note{la note précédente est écrite en oblique dans la marge gauche, à hauteur du corollaire}

\emph{Corollaire} Sous \struck{$X$ est une} ces conditions, (\ill{} \ill{} \ill{} \ill{} ultratopologiques) si \uncertain{$A$} est une partie \ill{} de $\mathcal{U}$, son adhérence pour $\tau'$, et \ill{} \ill{} \ill{}, est la \uncertain{même}. \struck{(\ill{} \ill{} \ill{} pour \ill{} de $T$.)}
\note{la lettre de « si $A$ est une partie » est écrite en surcharge ; la démonstration qui suit parle de $A$}

\struck{\ill{}} Soit \uncertain{en effet} $X \supset A$, fermé pour $\tau'$. Alors par les conditions 1°, 2°, 3°) plus \uncertain{haut},
\[
  X \supset \overline{X \cap \mathcal{U}}^{(\tau)} ,
\]
donc si $X \supset A$,
\[
  X \supset \overline{A}^{(\tau)} , \qquad \struck{\ill{}} \text{ en particulier}
\]
\[
  \overline{A}^{(\tau')} \supset \overline{A}^{(\tau)} , \qquad \text{OK.}
\]

Prouvons e). La première assertion (\uncertain{basée} sur le \uncertain{fait}) \struck{(les conditions (i) (ii) (iii) \uncertain{pour} $\tau'$ \ill{} \ill{} \ill{} \uncertain{connues})}\note{au-dessus du passage biffé : « \uncertain{considérées} \ill{} par $\tau'$ \ill{} \ill{} »} est \uncertain{évidente}. Montrons que la top. de $I$ est quotient pour $J \to I$, $I' \to I$. \struck{\uncertain{question} \ill{} de \ill{}} Cela signifie que \struck{$I'$, dans} pour une partie de $I'$ qui est \emph{$\widetilde{\tau}$-saturée} (i.e. de la forme $\varphi^{-1}(A)$), \ill{} celle-ci est fermée pour $\overline{\tau}'$ \uncertain{si} (\uncertain{et seulement si}) elle l'est pour $\widetilde{\tau}$, et sa trace sur $J'$ \uncertain{l'est} pour $\overline{\tau}$\note{« $\widetilde{\tau}$ » et « $\overline{\tau}$ » sont écrits sur deux lignes serrées en bas de page ; l'ordre de lecture est incertain}. Cela revient à \uncertain{ceci}, \struck{\ill{}}

\page{45}
\uncertain{revient} \struck{à} à l'\uncertain{interprétation} \uncertain{suivante}

\marginal{\ill{} \ill{} \uncertain{utile},}

\emph{Corollaire} \struck{Dans le cas des ultratopologies,}\note{au-dessus du passage biffé : « (Vu \uncertain{les} \uncertain{formules} de \uncertain{ventilation} \uncertain{en termes d'}ultratopologies) »} \struck{Soit} $A$ une partie de $T$ \struck{\ill{}} \struck{$\tau$.} Pour qu'elle soit fermée pour $\tau$, il faut et il suffit qu'elle soit a) fermée pour $\tau'$ et b) \uncertain{que} sa trace sur $T$ soit fermée pour $\tau$.
\note{la lettre lue $T$ dans « une partie de $T$ » ; la suite (« sa trace sur $T$ ») demande plutôt une partie de $S$}

[i.e. la top. $\tau$ sur $S$ est finale pour
\[
  S \xrightarrow{\ \mathrm{id}\ } S, \qquad T \xrightarrow{\ \mathrm{inc.}\ } S
\]
\struck{pour les topologies} $\tau'$ sur $S$, $\tau|T$ sur $T$.]

(C'est évidemment \struck{\ill{}} nécessaire.) Prouvons \ill{}

\note{le passage qui suit est barré de deux longs traits obliques ; il est transcrit tel qu'il se lit}
\struck{\ill{} \uncertain{Comme}} \uncertain{$X$} \ill{} fermé pour $\tau'$, (d'après (b)) \struck{donc} $\tau|\mathcal{U}$, \struck{\ill{}} et que $X \cap \mathcal{U}$ \add{$= X_{\mathcal{U}}$} est fermé pour $\tau|\mathcal{U}$, d'où $\overline{X_{\mathcal{U}}}^{(\tau)} = X_{\mathcal{U}}$ ; \uncertain{Comme} $X = \overline{X_{\mathcal{U}}}^{\tau} \cup X_B$ ; que $\overline{X_{\mathcal{U}}}^{(\tau)} \cap T \subset X_B$, et \uncertain{hypothèse}, \uncertain{il} \ill{} \uncertain{fermé} dans $T$ pour $\tau$ par hypothèse,

\uncertain{il} s'ensuit que $X$ \uncertain{est} \ill{} sur la formule
\[
  \overline{X}^{(\tau')} = \overline{X \cap \mathcal{U}}^{(\tau)} \cup \overline{X \cap T}^{(\tau'_0)}
\]
(\uncertain{cf.} p. 44, \uncertain{marge}), si $X$ est fermé pour $\tau'$, \uncertain{on a} \struck{$X_B$}, donc
\[
  X = \uncertain{\text{second membre}} = \overline{X \cap \mathcal{U}}^{(\tau)} \cup (X \cap T)
\]
donc fermé comme réunion de deux fermés.
\note{sous $(X \cap T)$, une accolade : « car $X \cap T$ fermé pour $\tau$, a fortiori pour $\tau'_0$ » ; la seconde topologie est lue $\uncertain{\tau'_0}$}

\page{46}
\emph{Corollaire} (Cas ultratopologiques)\note{la parenthèse est au-dessus de la ligne} Supposons (\ill{}), que la top. de $\mathcal{J}'$ soit quotient de celle de $\mathcal{J}$, \struck{Alors} i.e.
\[
  \varphi^* : \mathcal{J}' \to \mathcal{J} \text{ surjectif.}
\]
Alors $\varphi$ \ill{}.

\emph{Corollaire} (Ultratopologiques) Supposons $\varphi^*$ surjectif ; fibres \uncertain{0-connexes} — alors $\varphi$ itou. Supposons $\varphi$ \struck{\ill{}} homotopiquement propre i.e. les
\[
  \varphi^{-1}(J_{\leq x})
\]
\note{l'exposant de $\varphi$ est écrit $\ast$ surchargé en $-1$ ; sous $\leq$, un trait double, peut-être un $\not\leq$ corrigé}
« faiblement contractiles ». Alors kif-kif pour $\varphi^*$. \struck{(qui \ill{})} On est ramené à prouver que si $I$ admet un plus grand élément $x$ (\uncertain{donc} $x \notin J$, \uncertain{sinon} c'est f.c.)\note{« $x \notin J$ » est écrit au-dessus de la ligne}, alors $I'$ est f. contractile. Mais considérons l'unique \uncertain{élément} $x'$ \uncertain{qui s'envoie} sur $x$ ; \uncertain{je prétends} que c'est un plus grand élément. En effet, comme \ill{} plus grand \ill{} pour $\overline{\tau}'$, et pour $\overline{\tau}$, et $x' \in I \smallsetminus \mathcal{U}'$, \ill{} \ill{} \uncertain{c'est} \ill{}. OK !

\ill{} $=$

\emph{Corollaire} (Ultratopologiques)\note{la parenthèse est au-dessus de la ligne} Si $\varphi^{\ill{}}$ est homotopiquement propre à fibres \struck{\ill{}} faiblement contractiles (\uncertain{donc} une équivalence d'homotopie faible universelle), itou pour $\varphi$.

\page{47}
Mais je n'ai pas démontré \uncertain{entièrement} \ill{} (ou $\overline{\tau}'$) : pour prouver e), il faut montrer que $\tau'$ \uncertain{a} \ill{} les conditions dites.
\marginal{\uncertain{vrai ?} c'est faux !!}
\note{la note précédente est écrite en oblique dans la marge, pointant sur le mot souligné « \uncertain{Construction} » qui ouvre la troisième ligne}

Donc il faut voir que si $\tau''$, sur $S$, satisfait les conditions
\begin{itemize}
  \item[(i)] $\tau'' \leq \tau$
  \item[(ii)] $\tau''|T = \tau'_0$
  \item[(iii)] $\tau''|\mathcal{U} = \tau|\mathcal{U} \quad (\mathcal{U} = S \smallsetminus T)$
  \item[(iv)] Toute partie \struck{de} $X$ de $S$ \struck{telle} qui est fermée pour $\tau''$ et telle que $X \cap T$ fermée pour $\tau$, est fermée pour $\uncertain{\tau}$,
\end{itemize}
\marginal{i.e. top. finale pour $S_{(\tau'')} \xrightarrow{\mathrm{id}} S$, $T_{(\tau)} \xrightarrow{i} S$}
alors $\tau'' = \tau'$. On sait déjà, en \uncertain{vertu} de (i)\note{la référence est lue « (ii) » ou « (i) » ; seule (i) donne l'inégalité}, que $\tau'' \leq \tau'$, reste à voir l'inverse, i.e. que \struck{$X$ fermé pour $\tau''$} \uncertain{tout} fermé pour $\tau''$ est fermé pour $\tau'$, i.e. satisfait les conditions 1° 2° 3° de (b).
\note{« (b) » est entouré d'un cercle, au début de la ligne}
\begin{itemize}
  \item[1°)] $X \cap \mathcal{U}$ est fermé dans $\mathcal{U}$ pour $\tau|\mathcal{U}$ : OK puisque $\tau''|\mathcal{U} = \tau|\mathcal{U}$ (iii)
  \item[2°)] $X \cap T$ fermé dans $T$ pour $\tau'_0$ : OK puisque $\tau''|T = \tau'_0$ (ii)
  \item[3°)] $\overline{X_{\mathcal{U}}}^{(\tau)} \cap T \subset X \cap T$
\end{itemize}

Ceci n'est \uncertain{pas} \uncertain{démontré}, il faut utiliser (iv).

Il \uncertain{faudrait} \uncertain{prouver} que
\[
  \overline{X_{\mathcal{U}}}^{(\tau)} \subset \overline{X_{\mathcal{U}}}^{(\tau'')} \qquad \left(\subset \overline{X}^{(\tau'')} = X\right)
\]
\note{sous le dernier $X$ : « car $X$ fermé pour $\tau''$ »}
\uncertain{donc} $=$ ; et \uncertain{ça} \uncertain{donnerait} \uncertain{le} \uncertain{résultat}.

Donc il faudrait prouver que pour $A \subset \mathcal{U}$,
\[
  \overline{A}^{(\tau'')} = \overline{A}^{(\tau)} , \quad \text{i.e.} \quad \overline{A}^{(\tau'')} \text{ est } \tau\text{-fermé}
\]
\note{les derniers mots sont doublement soulignés}

\page{48}
et par la condition (iv), il reste à vérifier que
\[
  \overline{A}^{(\tau'')} \cap T \text{ est } \tau \text{ fermé} \qquad ??
\]
\marginal{C'est \uncertain{peut-être} faux ?}

C'est ultra faux \ill{}

\note{croquis : deux ovales l'un au-dessus de l'autre, reliés par une flèche verticale $\varphi$ (surchargée d'un signe biffé). L'ovale du haut, surmonté d'une accolade marquée $I'$, contient trois points : un point isolé en haut à gauche, relié au point $b$ par un trait épais hachuré, et deux points $a \to b$ reliés par une flèche. L'ovale du bas, marqué $J'$ à gauche et souligné d'une accolade marquée $I$, contient deux points reliés par une flèche. À gauche de l'ovale du haut, un gribouillis biffé}

$I'$ a trois éléments, $I$ deux. $J$ un seul, $J'$ deux. On prend $A = \{b\}$, son adhérence pour $\tau$ est $I'$ tout entier, pour $\tau'$ c'est $\{a, b\}$.

\note{un trait horizontal traverse la page ; le paragraphe suivant est marqué d'un double trait vertical dans la marge}

\uncertain{Finalement}, la condition \uncertain{ici} \uncertain{demandée} : \uncertain{imposer} à $\tau'$ \uncertain{pour} la \uncertain{construire} \uncertain{dans} $\underline{\underline{\mathcal{C}}}$, c'est la condition \uncertain{évidente} : pour $A \subset \mathcal{U} = S \smallsetminus T$, son adhérence pour $\tau$, et pour $\tau'$, \uncertain{est} la même (ce qui implique déjà $\tau|\mathcal{U} = \tau'|\mathcal{U}$). Cela donne l'adhérence \ill{} \ill{} d'une partie quelconque $A$, \uncertain{comme}
\[
  \overline{A}^{(\tau')} = \overline{A \cap \mathcal{U}}^{(\tau)} \cap \overline{(A \cap T)}^{(\tau'_0)}
\]
\note{le signe entre les deux adhérences est lu $\cap$ ; la même formule, en marge de la page 44, porte $\cup$}

\end{document}
